% Standalone section fragment for integration into the parent report.
% Sources to cite: eth_joint_cauchy, cor:smith;
% subexponential_sat_kernel_improvements, thm:domination.
\section{An inventory barrier that survives adaptive integral separation}
\label{sec:adaptive-inventory-barrier}

This result combines the exact full Cauchy Smith profile already proved~\cite{joint} in
\texttt{eth\_joint\_cauchy} with the pointwise optimality theorem~\cite{kernel} in
\texttt{subexponential\_sat\_kernel\_improvements}.  The added conclusion
also covers the sizes of explicit exact-weight tables.  It is a lower bound
on specified representations, not a lower bound for SAT or arbitrary
contraction algorithms.

Let $\mathcal O$ be the ring of integers of an unramified finite extension
of $\mathbb Q_2$, and put $R_P=\mathcal O/2^P\mathcal O$.
For an odd unit $\tau$, the integral Cauchy block
\[
 C_{jk}=\frac{2^\delta}{\tau+2(j-k)},\qquad 0\le j,k<J,
\]
has Smith valuations
\[
 e_j=\delta+4j-2s_2(j),
\]
truncated when they reach $P$.  Here $s_2(j)$ is binary popcount.
One proof applies the Cauchy determinant formula: an $h$-minor has
valuation
\[
 h\delta+h(h-1)+
 v_2\!\prod_{j<k}(x_k-x_j)+v_2\!\prod_{j<k}(y_k-y_j).
\]
Every Vandermonde on $h$ distinct integer points is divisible by
$\prod_{j=0}^{h-1}j!$, and consecutive indices attain equality.  The
successive determinantal-ideal valuations yield the displayed profile.
The full actual crossing orientation at valuation level $i$ has the same
profile with multiplicity $2^i$ and
$\delta=a+1-i$; this follows from its integral triangular factorization
with unit diagonal.  This latter fact is imported from the cited packet,
not claimed as new here.

\begin{lemma}[Integral separated profiles are dominated by Smith profiles]
\label{lem:profile-barrier}
Suppose a matrix over $R_P$ has its nonzero Smith valuations sorted as
$e_1\le\cdots\le e_t$, and admits an integral separated expression
\[
 K=\sum_{j=1}^{N}2^{w_j}a_jb_j^{\mathsf T},\qquad
 0\le w_1\le\cdots\le w_N<P.
\]
Then $N\ge t$ and $w_h\le e_h$ for $1\le h\le t$.
\end{lemma}
\begin{proof}
For each $1\le q\le P$, reduce modulo $2^q$.  The image of $K$ has
exactly $\#\{h:e_h<q\}$ nonzero cyclic summands, hence that many minimal
generators.  It is contained in the module generated by the
$\#\{j:w_j<q\}$ surviving target columns.  Over a finite chain ring,
the minimum number of generators equals the residue-field dimension of
the socle, and inclusion induces inclusion of socles.  Consequently
\[
 \#\{h:e_h<q\}\le\#\{j:w_j<q\}.
\]
The case $q=P$ gives $N\ge t$.  If $w_h>e_h$, the choice
$q=e_h+1$ contradicts the last inequality.  This is the optimality
argument of the incoming adaptive-separation theorem.
\end{proof}

\begin{theorem}[Explicit binary inventory lower bound]
\label{thm:adaptive-inventory}
Let $A,D,J,\delta,P$ be positive integers, let $u\ge1$ be real, and put
$M=AD$.  Suppose
\[
 D\ge16\sqrt u,\qquad
 \delta\le\frac D{\sqrt u},\qquad
 J\ge\left\lfloor\frac D{8\sqrt u}\right\rfloor,
 \qquad
 P\ge\frac{AD^2}{4u}.
\]
For the direct sum of $A$ copies of the preceding Cauchy block, every
integral separated expression with guaranteed weights $w_j$ has at least
\[
 2^{M/(16\sqrt u)}
\]
formal binary mode tuples satisfying both
$\sum_jw_j\beta_j<P$ and $\sum_j\beta_j\le M$.
\end{theorem}
\begin{proof}
Set $q=\lfloor D/(8\sqrt u)\rfloor$.  Since $D/(8\sqrt u)\ge2$,
\[
 q\ge\frac D{16\sqrt u}.
\]
The first $Aq$ Smith weights consist of the first $q$ weights in each
copy.  Their sum is bounded by
\begin{align*}
 A\sum_{j=0}^{q-1}(\delta+4j-2s_2(j))
 &\le A\bigl(q\delta+2q(q-1)\bigr)\\
 &\le\frac{5AD^2}{32u}<P.
\end{align*}
In particular these Smith entries are all nonzero at precision $P$.
Lemma~\ref{lem:profile-barrier} shows that the first $Aq$ weights of
any integral separated expression have no larger sum.  Every one of
their $2^{Aq}$ binary subsets is therefore feasible.  Its ordinary
degree is at most $Aq\le AD=M$.  The lower bound for $q$ completes the
proof.
\end{proof}

\begin{corollary}[An exact-weight list and a dense table must be large]
\label{cor:dense-table-barrier}
Under the hypotheses of Theorem~\ref{thm:adaptive-inventory}, at least one
integer $b$ with $0\le b<P$ has an exact-weight tuple list of size at least
\[
 \frac{2^{M/(16\sqrt u)}}P.
\]
An algorithm that explicitly materializes a group-function table with one
cell for every member of this list must write at least that many cells.
\end{corollary}
\begin{proof}
Partition the binary tuples from the theorem according to their integer
total weight.  There are at most $P$ possible weights.  Pigeonholing gives
the first claim, and explicit materialization gives the second.
\end{proof}

For a fixed number of groups, some directed group pair carries a fixed
positive fraction of any family of filled crossing blocks.  Applying the
corollary to that pair gives $2^{\Omega(M/\sqrt u)}/P$ formal indices.
In the reviewed schedules, $P=s^{O(1)}$, $u=O(\log s)$, the active wire
count is at least $s/\operatorname{polylog}s$, and the buffer ensures
$D/\sqrt u\gg\delta$.  Thus adaptive integral basis selection,
valuation-only pruning, the global ordinary-degree cap, and faster
matrix multiplication on fully materialized exact-weight tables cannot
by themselves remove this representation's square-root logarithmic
scale.  At the root $M=\Theta(s)$, it is
$2^{\Omega(s/\sqrt{\log s})}$.

When each full packet orientation is separated individually, the argument
also respects its row-degree cap: selecting $q$ modes in each of its
$R=2^i$ copies uses $Rq\le L/2$ modes, where $L=RD$.  If bases mix
different packets, only the asserted global degree cap follows without
additional support information.

No nonzero-entry lower bound is being asserted.  Distinct formal tuples
may cancel, yield zero, or have a compact joint representation.  The
bound does not constrain an algorithm that proves such cancellations,
compresses the group tables, bypasses their explicit construction, or
uses a different deformation.  It also does not prove that every SAT
encoding realizes the displayed filled family.  These are precisely the
ways a structural improvement could evade this local barrier.
